%% ==========================================================================
%%  FINAL REPORT --- SECTION 8, RELATED WORK
%%  Self-contained restatement of the survey's systems_both.tex, whose text
%%  refers to "this survey" and to survey-only sections. Quotations are the
%%  survey's, from papers read in full (reading tier A in references.bib).
%% ==========================================================================
\section{Related work}
\label{sec:related}

Our literature survey, submitted as the project's first milestone, covers the
field in full. This section places \sys{} against the work it builds on and the
work closest to it.

\paragraph{\pirsonaplain{}.} \pirsona{} is the system closest in spirit, and the
only one we found that spans collection, training and delivery. Its authors
describe private information retrieval and collaborative filtering as
``seemingly antithetical primitives'': PIR exists to hide which item a user
fetches, and collaborative filtering works by relating users through exactly
that. Its resolution is the harvest we reuse. Delivery uses the multi-server PIR
of Hafiz and Henry~\cite{hafiz2019}, the servers extract secret-shared consumption
from the query traffic itself, and training runs a four-party Boolean matrix
factorisation over it. \sys{} keeps the harvest and the PIR layer and replaces
the training core with \nudgeplain{}'s, one fewer non-colluding party and an
arithmetic rather than Boolean computation.

\paragraph{\nudgeplain{}.} \nudge{} is our training core. Users secret-share
ratings, three servers factorise by power iteration, and the servers compute each
user's scores in shared form. Its non-goal is our starting point: it ``relies on
other means (e.g., Apple's private relay, Tor, or cryptographic private
information retrieval) to let users fetch data items in a private way''. We add
three things it does not report: a private fetch composed with its training, a
measurement of what leaks when the fetch is not private, and a measurement of its
own differential-privacy extension against that leak.

\paragraph{Private training more broadly.} The first practical private matrix
factorisation, that of Nikolaenko et al.~\cite{nikolaenko2013}, evaluates
gradient descent inside a garbled circuit across two
servers.\footnote{We could not obtain this paper. The characterisation follows
\nudgeplain{}'s description of it, as our survey's did.} A garbled circuit's
communication scales with the circuit, and so with the whole rating matrix, and
later work moved to secret sharing. General frameworks such as MP-SPDZ~\cite{mpspdz2020} and the ABY
family~\cite{aby2015,aby22021} could run the same factorisation, and a generic
MPC baseline was part of our original evaluation plan. We did not build it, for
reasons Section~\ref{sec:limits} gives. The fixed-point costs that dominate our
measurements are a known theme: SecureML~\cite{secureml2017} and Escudero et
al.~\cite{escudero2020} both treat truncation as a first-class cost.

\paragraph{Function secret sharing and PIR.} The DPF is that of Boyle, Gilboa and
Ishai~\cite{gilboa2014,bgi2015,bgi2016}, and the comparison gate follows the
FSS-based secure computation of Boyle et al.~\cite{boyle2021}, which evaluates
comparison gates in a constant number of rounds by masking the input. Two-server
PIR from a DPF is a special case of Hafiz and Henry's family~\cite{hafiz2019}.
Single-server PIR such as SimplePIR~\cite{simplepir2023} removes the
non-collusion assumption at the price of a large client-side hint, and
Floram~\cite{floram2017} and \duoramplain{}~\cite{duoram2023} build distributed
oblivious memory from the same DPF machinery we use. \textsc{Sabre}'s
audit~\cite{sabre2022}, which proves that a DPF key encodes a genuine point
function, is the mechanism our malicious-client analysis names as the missing
piece.

\paragraph{Private retrieval without a recommender.} A separate line of work
retrieves privately over embeddings, from secure nearest-neighbour
search~\cite{sanns2020} to private web search~\cite{tiptoe2023} and more recent
sublinear schemes~\cite{pacmann2025,compass2025,panther2025,wally2024}. These
take a corpus and its embeddings as given and are not connected to a privately
trained model. \pirsonaplain{} predates almost all of them, which is part of why
the composition we build had not been attempted with a modern training core.

\paragraph{What leaks through the output.} Calandrino et
al.~\cite{calandrino2011} address the privacy risks of collaborative filtering
itself.\footnote{We could not obtain this paper either, and cite it only for the
claim its title makes.} Our reconstruction attack concerns a narrower channel. It
needs no access to the recommender's output, only to which items a user fetched,
together with the public $\itemembed$ that \nudgeplain{}'s design releases.
